\abstract{
    Cedar is a protoype programming language and compiler that introduces
    declarative byte-level data-layout specication to the C programming
    ecosystem.
    Low-level software such as device drivers, network stacks, and embedded
    systems often require precise control over how data structures are
    represented in memory.
    In the C programming langauge, these details are typically implicit,
    expressed through field ordering, alignment directives, or
    implementation-sepcific pragmas. Such mechanisms are powerful
    but error-prone, and layout inconsistencies are a major source of system
    failures and security vulenerabilities.

    Cedar builds on ideas from Dargent, a verified layout description
    language integrated into the Cogent programming langauge.
    Cedar also borrows from CompCert, the formally verified C compier.
    Rather than pursuing full formal verification, Cedar explores how
    Dargent's declarative approach can be realized directly within
    conventional C development. The language allows programmers to
    specify the structure and placement of data fields using high-level
    constructs such as records, absolute and relative byte offsets,
    endianness annotations, and fixed-length arrayss. Each layout
    is compiled into C code that represents the structure as a flat array
    of 32-bit words, with automatically generated accessor and mutator
    functions for every field.

    The implementation comprises a four-stage pipeline: parsing, lowering,
    checking, and code generation.
    This mirrors the refinement structure of CompCert while remaining small and
    inspectable.
    The prototype supports nested records, per-field endianness, and relative
    or absolute placements, and it generates portable C deo output that has
    been tested to compile succesfully with both GCC and CompCert.
    Evaluation on representative layout examples demonstrate that Cedar can
    express complex C-style structures clearly and that its generated code
    matches the offsets and sizes produced by standard compilers.

    By making layout intent explicit and checkable, Cedar offers a practical
    step towrads trustworthy data-layout specification without imposing a 
    new programming model or sacrificing compatibility with exisiting toolchains.
    The project highlights a path between unverified low-level programming
    and fully verified systems: a middle ground where clarity and correctness
    can coexists with C's performance and flexibility.
}